% Generated by scripts/assemble.py; corrections belong in source/corrections.json.
\examhead{156}{华南理工大学 2025 年数学分析试卷}{149}
% p150-b01; source PDF page 150
\blockanchor{p150-b01}
\question{1}(12 分) 已知数列 $x_0 = a > 0$, $x_n = \arctan x_{n-1} (n = 1, 2, \cdots)$. 证明:\par
% p150-b02; source PDF page 150
\blockanchor{p150-b02}
(1) $\lim_{n\to \infty}x_n = 0.$\par
% p150-b03; source PDF page 150
\blockanchor{p150-b03}
(2) $\lim_{n\to \infty}\sqrt{\frac{2n}{3}} x_n = 1.$\par
% p150-b04; source PDF page 150
\blockanchor{p150-b04}
\question{2}(12 分) 证明函数 $f(x) = \begin{cases} \sqrt{x} \ln x, & x > 0; \\ 0, & x = 0. \end{cases}$ 在 $[0, +\infty)$ 上一致连续.\par
% p150-b05; source PDF page 150
\blockanchor{p150-b05}
\question{3}(12 分) 证明不等式 $\frac{1}{\sin^2 x} \leq \frac{1}{x^2} + 1 - \frac{4}{\pi^2} \left(0 < x \leq \frac{\pi}{2}\right)$.\par
% p150-b06; source PDF page 150
\blockanchor{p150-b06}
\question{4}(14 分) 设 $f(x, y) = \begin{cases} g(x, y) \sin \frac{1}{\sqrt{x^2 + y^2}}, & (x, y) \neq (0, 0); \\ 0, & (x, y) = (0, 0). \end{cases}$ 证明:\par
% p150-b07; source PDF page 150
\blockanchor{p150-b07}
(1) 若 $g(0,0) = 0$, $g(x,y)$ 在 $(0,0)$ 处可微且 $\mathrm{d}g(0,0) = 0$. 则 $f(x,y)$ 在 $(0,0)$ 处可微, 且 $\mathrm{d}f(0,0) = 0$.\par
% p150-b08; source PDF page 150
\blockanchor{p150-b08}
(2) 若 $g(x, y)$ 在 $(0, 0)$ 处有偏导数且 $f(x, y)$ 在 $(0, 0)$ 处可微, 则 $\mathrm{d}f(0, 0) = 0$.\par
% p150-b09; source PDF page 150
\blockanchor{p150-b09}
\question{5}(12 分) 设 $f(x)$ 在 $[0,1]$ 上连续, 在 $(0,1)$ 内可导, $f(0) = 1$. 证明:\par
% p150-b10; source PDF page 150
\blockanchor{p150-b10}
(1) 存在 $\xi \in (0,1)$, 使得 $f(1) = (1 + \xi)f'(\xi)\ln 2 + 1$.\par
% p150-b11; source PDF page 150
\blockanchor{p150-b11}
(2) $\lim_{n\to \infty}n(\sqrt[n]{1 + x} -1) = \ln (1 + x),x\in (0,1).$\par
% p150-b12; source PDF page 150
\par\Needspace{6\baselineskip}
\blockanchor{p150-b12}
\question{6}(12 分) 设 $f(x)$ 在 $[0,1]$ 上可积, 证明: 存在折线函数列 $\{f_n(x)\}_{n=1}^{\infty}$, 使得\par
\nopagebreak[4]
% p150-b13; source PDF page 150
\blockanchor{p150-b13}
\displayformula{\int_ {0} ^ {1} f (x) \mathrm{d} x = \lim _ {n \rightarrow \infty} \int_ {0} ^ {1} f _ {n} (x) \mathrm{d} x.}
% p150-b14; source PDF page 150
\blockanchor{p150-b14}
\question{7}(12 分) 求幂级数 $\sum_{n=1}^{\infty} \frac{x^{4n-1}}{4n+1}$ 的收敛域与和函数 $S(x)$.\par
% p150-b15; source PDF page 150
\par\Needspace{6\baselineskip}
\blockanchor{p150-b15}
\question{8}(12 分) 计算曲面积分 $F(t) = \iint_{x + y + z = t} f(x, y, z) \, \mathrm{d}S$, 其中\par
\nopagebreak[4]
% p150-b16; source PDF page 150
\blockanchor{p150-b16}
\displayformula{f (x, y, z) = \left\{ \begin{array}{l l} 1 - x ^ {2} - y ^ {2} - z ^ {2}, & x ^ {2} + y ^ {2} + z ^ {2} \leq 1; \\ 0, & x ^ {2} + y ^ {2} + z ^ {2} > 1. \end{array} \right.}
% p150-b17; source PDF page 150
\blockanchor{p150-b17}
\question{9}(12 分) 设 $\int_{0}^{+\infty} f(x) \, \mathrm{d}x$ 收敛, 且 $\lim_{x \to +\infty} f(x)$ 有意义(有限数或 $+\infty$ 或 $-\infty$). 证明: $\lim_{x \to +\infty} f(x) = 0$.\par
% p150-b18; source PDF page 150
\par\Needspace{6\baselineskip}
\blockanchor{p150-b18}
\question{10}(12 分) 设 $f(x)$ 在 $[0, +\infty)$ 上连续且 $\lim_{x \to +\infty} f(x) = k$, 试证:\par
\nopagebreak[4]
% p150-b19; source PDF page 150
\blockanchor{p150-b19}
\displayformula{\int_ {0} ^ {+ \infty} \frac {f (a x) - f (b x)}{x} \mathrm{d} x = [ f (0) - k ] \ln \frac {b}{a} (a > 0, b > 0).}
% p150-b20; source PDF page 150
\par\Needspace{6\baselineskip}
\blockanchor{p150-b20}
\question{11}（13分）设 $u = x + y, v = \frac{1}{x} + \frac{1}{y}$，试用新变量 $u, v$ 变换等式\par
\nopagebreak[4]
% p150-b21; source PDF page 150
\blockanchor{p150-b21}
\displayformula{x ^ {2} \frac {\partial^ {2} z}{\partial x ^ {2}} - (x ^ {2} + y ^ {2}) \frac {\partial^ {2} z}{\partial x \partial y} + y ^ {2} \frac {\partial^ {2} z}{\partial y ^ {2}} = 0.}
% p150-b22; source PDF page 150
\blockanchor{p150-b22}
(假设所有出现的二阶偏导数都连续)\par
% p150-b23; source PDF page 150
\blockanchor{p150-b23}
\question{12}(15 分) 设函数列 $f_{n}(x) = nxe^{-n^{2}x^{2}}$, $x \in [0,1]$, $n = 1,2,\dots$. 证明:\par
% p151-b00; source PDF page 151
\blockanchor{p151-b00}
(1) $\{f_n(x)\}_{n=1}^{\infty}$ 在 $[0,1]$ 上收敛于 $f(x) \equiv 0$.\par
% p151-b01; source PDF page 151
\blockanchor{p151-b01}
(2) $\{f_n(x)\}_{n=1}^{\infty}$ 在 $[0,1]$ 上是否一致收敛于 $f(x) \equiv 0$? 判断并给出理由.\par
% p151-b02; source PDF page 151
\blockanchor{p151-b02}
(3) $\{f_n(x)\}_{n=1}^{\infty}$ 在 $[0,1]$ 上积分平均收敛于 $f(x) \equiv 0$, 即 $\lim_{n \to \infty} \int_0^1 |f_n(x) - f(x)| \, \mathrm{d}x = 0$.\par
